%! AP Statistics — Unit 7, Lesson 7.3 — Guided Notes
\documentclass[10pt]{article}
\usepackage{apstats-article}
\usepackage{apstats-boxes}
\newcommand{\TallMath}[1]{$\displaystyle #1\rule[-1.4em]{0pt}{3.2em}$}

\begin{document}

\pageheader{Unit 7, Lesson 7.3}{Guided Notes}
\namedateperiod

\begin{objectivebox}
By the end of this lesson, I will be able to:
\begin{itemize}
    \item[$\square$] \writeline
    \item[$\square$] \writeline
    \item[$\square$] \writeline
    \item[$\square$] \writeline
\end{itemize}
\end{objectivebox}

\begin{vocabbox}
\termblanklong{Interpretation of a confidence interval}
\termblanklong{``$C$\% confident'' (what it really means)}
\termblanklong{Width of a confidence interval}
\termblanklong{Justifying a claim using a CI}
\end{vocabbox}

\begin{hookbox}
From Lesson 7.2: \textbf{95\% CI for mean sodium intake: $(2615, 3065)$ mg.}

Below are two interpretations. Which is correct? What is wrong with the other?

\begin{itemize}
    \item[A.] ``There is a 95\% probability that the true mean is between 2615 and 3065 mg.''
    \item[B.] ``We are 95\% confident that the interval (2615, 3065) mg captures
              the true mean daily sodium intake for all college students.''
\end{itemize}

Correct: \blank{2cm} \quad Why is the other wrong? \writelines{2}
\end{hookbox}

\begin{notesbox}{Correct Interpretation Template (UNC-4.S)}
``We are \blank{2cm}\% confident that the interval
$(\rule{1.5cm}{0.4pt},\ \rule{1.5cm}{0.4pt})$ \blank{2cm}
captures the true mean \blank{4cm} for \blank{4cm}.''

\vspace{0.1in}
\textbf{What does ``$C$\% confident'' actually mean?} (UNC-4.S.1)\\
In \blank{8cm} with the same $n$, approximately \blank{2cm}\% of
all confidence intervals constructed this way will capture the true mean $\mu$.

\vspace{0.1in}
\textbf{Three common errors to avoid:}
\begin{enumerate}
    \item Using ``\blank{4cm}'' instead of ``confident'' (implies $\mu$ is random).
    \item Failing to reference the \blank{4cm} or \blank{4cm} (UNC-4.S.3).
    \item Claiming the interval contains \blank{4cm}\% of the data values (wrong ---
          it estimates the \blank{4cm}, not individual values).
\end{enumerate}
\end{notesbox}

\begin{notesbox}{Width of a Confidence Interval (UNC-4.U)}
Width $= 2 \times ME = 2 \cdot t^* \cdot \dfrac{s}{\sqrt{n}}$

\vspace{0.1in}
\begin{tabular}{lll}
    \textbf{Change} & \textbf{Effect on Width} & \textbf{Why} \\
    \hline
    Increase $n$ & Width \blank{3cm} & $ME \propto 1/\sqrt{n}$; to halve width, \blank{4cm} $n$ \\
    Increase $C$ & Width \blank{3cm} & Larger $t^*$ is needed \\
    Increase $s$ & Width \blank{3cm} & More variability in the data \\
\end{tabular}

\vspace{0.1in}
Width is proportional to $1/\sqrt{n}$: quadrupling $n$ \blank{4cm} the width (UNC-4.U.2).
\end{notesbox}

\begin{notesbox}{Justifying a Claim Using a CI (UNC-4.T)}
Given a confidence interval $(\ell, u)$ and a claimed value $\mu_0$:

\vspace{0.1in}
\begin{tabular}{ll}
    \textbf{If $\mu_0$ is outside the CI} & \blank{6cm} evidence that
          $\mu \ne \mu_0$ \\
    \textbf{If $\mu_0$ is inside the CI} & \blank{6cm} evidence against the claim \\
\end{tabular}

\vspace{0.1in}
\textbf{Example:} Sodium CI $= (2615, 3065)$ mg. Claimed value: $\mu_0 = 2300$ mg.

Is 2300 inside or outside the interval? \blank{3cm}

Conclusion: \writelines{2}

\vspace{0.1in}
\textbf{One-sided claims:} If the claim is ``$\mu > 2500$,'' check whether
\blank{6cm} of the CI lies above 2500.
\end{notesbox}

\begin{practicebox}
\textbf{Guided Practice:} A 90\% CI for the mean wait time at a clinic is
$(18.3, 24.7)$ minutes. The clinic manager claims the mean wait time is 20 minutes.

\begin{enumerate}
  \item Write a correct interpretation of the interval. \writelines{2}
  \item Does the interval provide convincing evidence against the manager's claim?
        Explain. \writelines{2}
  \item If the sample size were quadrupled (same $\bar x$ and $s$), describe the
        effect on the margin of error and the interval width. \writelines{2}
\end{enumerate}
\end{practicebox}

\end{document}
